\section{Local harmonic analysis in the theta estimates}
\label{sec:theta-local-audit}

This section reconstructs several analytic and local arithmetic steps that
enter the proposed estimates for the theta maps.  Every lemma proved here is
independent of the Riemann hypothesis.  The final estimate for the actual
global family is stated with its remaining arithmetic hypotheses displayed.
In particular, a local stationary phase formula is not a proof of the global
estimate for the sum of the orbital terms.

The source uses the normalized operators
\begin{equation}\label{th-normalized-K}
 K_{uv}=(c_uc_v)^{-1/2}P_u^*P_v,
 \qquad
 c_u=\frac{\kappa\,\operatorname{vol}(N_u)
                   \operatorname{vol}(H_u)}{\dim\pi_u}.
\end{equation}
Here $P_u:C_u\longrightarrow\mathcal H$ is a Poincar\'e-series map,
$C_u$ is a finite-dimensional Hilbert space, $\pi_u$ is a nonzero compact
stabilizer representation, and $\kappa$ is the volume of the finite-level
subgroup.  The source incorporates $\kappa$ into its symbol
$\operatorname{vol}(N_u)$.  All inner products below are conjugate linear
in the first variable.  Thus $K_{uv}:C_v\longrightarrow C_u$.

\subsection{The local rings and measures actually required}

Let $K$ be a finite product of finite extensions of $\mathbb Q_p$.  A
compact open additive subgroup of $K^n$ will be called a lattice.  For
$a\in K^\times$, write
\[
 |a|_K=\prod_{K_j\mid K}|a_j|_{K_j},
\]
where the absolute value in a field factor gives a uniformizer the value
$(\#k_j)^{-1}$.  For a matrix $B\in\operatorname{GL}_n(K)$,
linear change of variables multiplies additive Haar measure by
$|\det B|_K$.

Fix a continuous additive character $\psi:K\to\mathbb C^\times$ that is
nontrivial on every field factor.  Choose the self-dual Haar measure, and
use its product measure on $K^n$.  For a lattice $L\subset K^n$, put
\[
 L^\perp=\{y\in K^n:\psi(x^ty)=1\text{ for all }x\in L\}.
\]
Then $(L^\perp)^\perp=L$ and
\begin{equation}\label{th-dual-vol}
 \operatorname{vol}(L)\operatorname{vol}(L^\perp)=1,
 \qquad
 \int_L\psi(x^ty)\,dx=
 \operatorname{vol}(L)\,1_{L^\perp}(y).
\end{equation}
For completeness, the second identity follows by translating the integral
by an element $x_0\in L$: if the character is nontrivial, that translation
multiplies the integral by $\psi(x_0^ty)\ne1$.  The duality and volume
identities reduce by a linear change of basis to fractional ideals in a
local field.  If the annihilator of its ring of integers $O$ is $c^{-1}O$,
the measure for the character $x\mapsto\psi_0(cx)$ is
$|c|^{1/2}\,dx_0$, where $\psi_0$ has annihilator $O$ and
$\operatorname{vol}_{dx_0}(O)=1$.  The identities for fractional ideals
follow immediately.  A product of these calculations proves
\eqref{th-dual-vol}; a general additive lattice is a free
$\mathbb Z_p$-module, so the same calculation can alternatively be made
over $\mathbb Q_p$ using the trace pairing.

The two local orders in the source are as follows.

\begin{enumerate}
\item Away from the two distinguished split primes, $K$ is an unramified
quadratic algebra over a local field $F$, $R=O_K$, and the specified
character has annihilator $O_K$.  This includes the algebra $F\times F$.
The chosen measure has $\operatorname{vol}(R)=1$.
\item At a distinguished split prime, write $K=F\times F$, let $\varpi$
be a uniformizer of $F$, and let $\psi_F$ have annihilator $O_F$.  The
order and character are
\begin{equation}\label{th-split-order}
 R_\varpi=\{(x,y)\in O_F^2:x-y\in\varpi O_F\},
 \qquad
 \psi_R(x,y)=\psi_F((x-y)/\varpi).
\end{equation}
\end{enumerate}

\begin{lemma}\label{th-split-self-dual}
The order $R_\varpi$ is its own annihilator for $\psi_R$.
If $q=\#(O_F/\varpi O_F)$ and product Haar measure gives $O_F^2$
volume one, self-dual Haar measure for $\psi_R$ is $q$ times that
product measure.  In particular, it gives $R_\varpi$ volume one.
\end{lemma}

\begin{proof}
Write an element of $R_\varpi$ as $(y+\varpi z,y)$, with $y,z\in O_F$.
An element $(a,b)\in F^2$ annihilates this order if and only if
\[
 \psi_F\bigl(az+(a-b)y/\varpi\bigr)=1
 \quad(y,z\in O_F).
\]
This is equivalent to $a\in O_F$ and $(a-b)/\varpi\in O_F$, which
is precisely $(a,b)\in R_\varpi$.  The characters on the two factors
have scale factors $\varpi^{-1}$ and $-\varpi^{-1}$.  Their self-dual
measures therefore contribute $q^{1/2}$ each.  Finally
$[O_F^2:R_\varpi]=q$.
\end{proof}

Consequently it would be incorrect to replace all local orders by maximal
orders while retaining the same Haar measure.  The source does retain the
normalization in Lemma~\ref{th-split-self-dual}.  The proofs below allow
this nonmaximal order; they use the additive pairing, rather than a claim
that $R$ is a discrete valuation ring.

\subsection{A Gaussian integral with all normalizing factors}

Let $B\in\operatorname{GL}_n(K)$ be symmetric and put
$q_B(x)=\tfrac12x^tBx$.  Characteristic zero is essential here.
For a lattice $L\subset K^n$, impose the explicit condition
\begin{equation}\label{th-gauss-admissible}
 \psi\bigl(q_{B^{-1}}(\xi)\bigr)=1\qquad(\xi\in L^\perp).
\end{equation}
Such lattices exist: take $L$ sufficiently large, so that $L^\perp$
is contained in a neighborhood on which the indicated continuous
character-valued function is one.

\begin{lemma}\label{th-gaussian}
If \eqref{th-gauss-admissible} holds, then
\[
 \left|\int_L\psi(q_B(x))\,dx\right|=|\det B|_K^{-1/2}.
\]
The integral is independent of the lattice $L$ satisfying that condition.
Accordingly there is a well-defined complex number $\gamma_\psi(B)$,
of absolute value one, for which
\begin{equation}\label{th-gaussian-value}
 \int_L\psi(q_B(x))\,dx
   =\gamma_\psi(B)|\det B|_K^{-1/2}.
\end{equation}
It is unchanged under congruence:
$\gamma_\psi(C^tBC)=\gamma_\psi(B)$ for $C\in\operatorname{GL}_n(K)$.
\end{lemma}

\begin{proof}
Set $U=B^{-1}L^\perp$.  Condition \eqref{th-gauss-admissible} says
$\psi(q_B(u))=1$ for $u\in U$.  Polarization gives
$\psi(u^tBv)=1$ for $u,v\in U$.  Since $B$ is symmetric,
$U^\perp=BL$.  It follows that $BU=L^\perp\subset BL$, or $U\subset L$.
Writing $I_L=\int_L\psi(q_B(x))\,dx$ and making the change of variables
$x=y+u$ gives
\begin{align*}
 |I_L|^2
 &=\int_L\psi(q_B(u))
       \int_L\psi(y^tBu)\,dy\,du\\
 &=\operatorname{vol}(L)
       \int_{L\cap B^{-1}L^\perp}\psi(q_B(u))\,du\\
 &=\operatorname{vol}(L)\operatorname{vol}(B^{-1}L^\perp)
   =|\det B|_K^{-1}.
\end{align*}
All these integrals are over compact sets, so Fubini causes no
convergence issue.

If $a\notin L$, duality supplies $\xi\in L^\perp$ with
$\psi(\xi^ta)\ne1$.  Put $u=B^{-1}\xi\in U\subset L$.
Translation by $u$ multiplies the integral of $\psi(q_B)$ over $a+L$
by $\psi(\xi^ta)$: the other factors are one because
$\psi(q_B(u))=1$ and $\xi\in L^\perp$.  That coset integral is zero.
Thus enlarging $L$ to any lattice leaves $I_L$ unchanged.  For two
admissible lattices, enlarge both to their sum.  This proves independence.
Finally substitute $x=Cy$ in the integral.  The factor $|\det C|_K$
is canceled by the change in $|\det(C^tBC)|_K^{1/2}$.
\end{proof}

For $m\ge1$, let $J_m$ denote the matrix all of whose entries are one,
and put
\[
 H_m(a)=a^{-1}(I_m+J_m),\qquad
 \gamma_m(a)=\gamma_\psi(H_m(a))\quad(a\in K^\times).
\]
The determinant lemma, or the decomposition into the line spanned by
$(1,\ldots,1)$ and its sum-zero complement, gives
\begin{equation}\label{th-hessian-det}
 \det H_m(a)=(m+1)a^{-m},\qquad
 (I_m+J_m)^{-1}=I_m-\frac{J_m}{m+1}.
\end{equation}
The identities remain valid when the residue characteristic divides
$m+1$: division takes place in a characteristic-zero field.  The
congruence assertion in Lemma~\ref{th-gaussian} gives
$\gamma_m(ab^2)=\gamma_m(a)$.  Thus
\begin{equation}\label{th-gauss-H}
 \int_L\psi\bigl(\tfrac12x^tH_m(a)x\bigr)\,dx
 =\gamma_m(a)|m+1|_K^{-1/2}|a|_K^{m/2}
\end{equation}
whenever \eqref{th-gauss-admissible} is satisfied.

\begin{remark}\label{th-conductor}
There is no omitted conductor factor in \eqref{th-gauss-H}, because the
Haar measure is self-dual for the displayed character.  If a field factor
uses $\psi_c(x)=\psi_0(cx)$ and one insists on the measure for $\psi_0$,
the right side must instead be multiplied by $|c|^{-m/2}$.
Moreover the admissibility condition changes with $c$.  One may not
change the character and keep both the measure and the smallness
threshold unchanged.
\end{remark}

\subsection{An elementary analytic change of variables}

We record the analytic fact needed for stationary phase, including a
proof that does not assume that $2$ or $m+1$ is a unit of $O_F$.

\begin{lemma}\label{th-morse}
Let $F$ be a finite extension of $\mathbb Q_p$ and let $f$ be analytic
near $0\in F^m$.  Suppose $df(0)=0$ and the symmetric Hessian
$H=D^2f(0)$ is invertible over $F$.  On a sufficiently small ball
$B=(\varpi^aO_F)^m$ there is an analytic bijection
$\Phi:B\to B$ such that
\[
 f(y)=f(0)+\tfrac12\Phi(y)^tH\Phi(y).
\]
The map preserves Haar measure, and for every $b\ge a$ and $y,y'\in B$,
\[
 \Phi(y)-\Phi(y')\in(\varpi^bO_F)^m
 \quad\Longleftrightarrow\quad y-y'\in(\varpi^bO_F)^m.
\]
For an analytic family the ball can be chosen uniformly under the
following precise hypothesis: on a common polydisc of radius $r>0$,
the convergent power series $\sum_\alpha c_\alpha y^\alpha$ have
uniformly bounded Gauss norms
$\sup_\alpha |c_\alpha|r^{|\alpha|}$, and the inverse Hessians have
uniformly bounded matrix norms.  Shrinking the common radius is allowed.
\end{lemma}

\begin{proof}
Expanding the convergent power series and assigning each monomial of
degree at least two to two of its variables writes
\[
 f(y)-f(0)=\tfrac12y^t(H+C(y))y,
 \qquad C(0)=0,
\]
with $C(y)$ symmetric and analytic.  One explicit assignment is to put a
monomial into a diagonal entry if some variable occurs at least twice,
and otherwise into the two symmetric entries corresponding to its
first two variables.  Only the fixed factors $2$ and $1/2$ are
introduced; after shrinking the polydisc the resulting series converge.
Put $A(y)=H^{-1}C(y)$.  The matrix $A(y)$ is self-adjoint for the
bilinear form $H$, namely $A(y)^tH=HA(y)$.

The formal series
\[
 S(y)=\sum_{j\ge0}\binom{1/2}{j}A(y)^j
\]
converges on a sufficiently small polydisc.  Indeed the $p$-adic
absolute values of the coefficients are bounded by an exponential in
$j$, whereas the matrix norm of $A(y)$ tends to zero with $y$.
It satisfies $S(y)^2=I+A(y)$ and $S(y)^tH=HS(y)$.  Therefore
$S(y)^tHS(y)=H+C(y)$.  Set $\Phi(y)=S(y)y$.

Here are the required invertibility and measure details.  Shrink $B$
until $g(y)=\Phi(y)-y$ has Lipschitz constant strictly less than one,
and $g(0)=0$.  This is possible because its convergent power series
has no constant or linear term.  If $b\in B$, the map
$y\mapsto b-g(y)$ is a strict contraction of the complete ball $B$,
so it has a unique fixed point.  This proves bijectivity.
Furthermore
\[
 \|\Phi(y)-\Phi(y')\|=\|y-y'\|
\]
in the supremum norm, since the error has strictly smaller norm.
The map consequently permutes the residue classes of each smaller
ball.  These classes generate the Borel sets of $B$ and all have
the same volume; hence $\Phi$ preserves Haar measure.  The power-series
bounds just used can be made uniform when the coefficients and
$H^{-1}$ satisfy the stated uniform bounds.  Indeed on any strictly
smaller radius the uniform Gauss bound gives uniform convergence of
the tails and of their divided difference estimates.
\end{proof}

\subsection{The stationary phase formula used in the local induction}

For $D\in K^\times$ define, on $(K^\times)^m$,
\[
 f_D(x_1,\ldots,x_m)=x_1+\cdots+x_m+
                         \frac{D}{x_1\cdots x_m}.
\]
Fix a compact subset $Q\subset K^\times$, a compact open set
$C\subset(K^\times)^m$, and a lattice $L_1\subset K^m$.
The amplitudes considered below are all locally constant functions
$A:K^m\to\mathbb C$, supported in $C$ and invariant under translation
by $L_1$.  No bound on their values is required.  Choose $L_1$ smaller
if necessary; doing so enlarges the allowed class of amplitudes.

\begin{theorem}\label{th-stationary}
There is a product neighborhood $V$ of zero in $K$, depending on
$K,\psi,m,Q,C,L_1$, such that for $t\in V\cap K^\times$ and $D\in Q$,
\begin{align}\label{th-stationary-formula}
 &\int_{(K^\times)^m}A(x)\psi(f_D(x)/t)\,dx\notag\\
 &\quad=|t|_K^{m/2}|m+1|_K^{-1/2}
   \sum_{\substack{z\in K^\times\\z^{m+1}=D}}
    \gamma_m(zt)|z|_K^{m/2}
    \psi((m+1)z/t)A(z,\ldots,z).
\end{align}
The sum is finite.  If $K$ is a field it has at most $m+1$ terms;
if $K$ has $d$ field factors it has at most $(m+1)^d$ terms.
\end{theorem}

\begin{proof}
We first take $K=F$ to be a field.  Direct differentiation gives
\[
 \frac{\partial f_D}{\partial x_i}
   =1-\frac{D}{x_i\prod_jx_j}.
\]
All derivatives vanish if and only if $x_1=\cdots=x_m=z$ and
$z^{m+1}=D$.  At that point the diagonal Hessian entries are $2/z$
and the off-diagonal entries are $1/z$.  Thus the Hessian is $H_m(z)$,
which is invertible by \eqref{th-hessian-det}, and
$f_D(z,\ldots,z)=(m+1)z$.

The set of possible roots $z$, for $D\in Q$, is compact in $F^\times$:
their absolute values lie in finitely many valuation classes, bounded
above and away from zero.  The coefficients of the Taylor expansion
at these roots, and the inverse Hessians, are therefore bounded
uniformly.  More precisely, choose a common expansion radius strictly
smaller than the least absolute value of these roots.  On that radius
the reciprocal factors in $f_D$ have geometrically convergent series
with uniformly bounded Gauss norms.  This is the uniform analytic
hypothesis of Lemma~\ref{th-morse}, which supplies a common sufficiently small
ball $B=(\varpi^aO_F)^m\subset L_1$.  We may also arrange that distinct
critical points have disjoint translates of $B$.  To see uniform
separation, otherwise there would be distinct roots $z_j,z'_j$ with
$z_j-z'_j\to0$; after taking a convergent subsequence, their common
limit $z$ is nonzero.  The function $w\mapsto w^{m+1}$ has derivative
$(m+1)z^m\ne0$, so is injective on a neighborhood of $z$, a
contradiction.  Its local injectivity follows by exactly the contraction
argument in Lemma~\ref{th-morse}, applied after its nonzero linear term
has been divided out.

On a critical ball, $A$ is constant.  Applying the measure-preserving
change of variables and then \eqref{th-gauss-H} gives
\[
 \int_{(z,\ldots,z)+B}A(x)\psi(f_D(x)/t)\,dx
 =A(z,\ldots,z)\psi((m+1)z/t)
   \gamma_m(zt)|m+1|^{-1/2}|zt|^{m/2}.
\]
The use of the Gaussian formula is justified uniformly for small $t$.
In fact
\[
 (H_m(z)/t)^{-1}
       =tz\left(I_m-\frac{J_m}{m+1}\right).
\]
On the fixed compact lattice $B^\perp$, the associated quadratic
form tends uniformly to zero as $t\to0$, with $z$ in its compact
set.  It is consequently in the kernel of $\psi$ for all sufficiently
small $t$.  This proves the admissibility condition, including at
primes dividing $2(m+1)$.

We give the uniform vanishing argument on the complement.  Near any
pair $(D_0,x_0)$ with $\nabla f_{D_0}(x_0)\ne0$, choose $j$ for which
$|a|=|\partial_jf_{D_0}(x_0)|$ is the largest absolute value among
the partial derivatives.  In particular $a\ne0$.  On sufficiently small product balls
in $D$ and $x$, the map
\[
 x\longmapsto(f_D(x),x_1,\ldots,\widehat{x_j},\ldots,x_m)
\]
is a bijection onto a product of balls.  Its first ball is a translate
of $a\varpi^bO_F$; the other balls have radius $|\varpi|^b$.
Its Jacobian has constant absolute value $|a|$.  These facts follow
by dividing the first coordinate by $a$ and applying the same strict
contraction argument: the nonlinear remainder and the variation in
the derivative can both be made smaller than the linear term.
For fixed other coordinates, this is a one-variable bijection of
balls, preserving their normalized Haar measures.  Shrink the
balls further so that $A$ is constant there.  Integration in the
first coordinate now vanishes as soon as
$\psi$ is nontrivial on $t^{-1}a\varpi^bO_F$.

Only finitely many such bounds are needed.  Here is the compactness
point needed for that assertion.  A root $z_0^{m+1}=D_0$ continues
uniquely, in a fixed small root ball, for all $D$ sufficiently near
$D_0$, by the preceding local inverse argument.  Roots outside those
balls cannot suddenly enter the fixed compact set under consideration:
a convergent subsequence would give an additional root at $D_0$.
Thus after shrinking a neighborhood of each $D_0$, the critical
balls have fixed disjoint locations and their centers vary within
those balls.  The remaining part of $C$ is a compact open set without
critical points, uniformly for these $D$.  It is a finite disjoint
union of balls of the kind just described.  A finite cover of $Q$
provides finitely many choices of $a,b$.  Choose $V$ small enough to
make all corresponding character integrals vanish.  Summing the
critical-ball integrals proves \eqref{th-stationary-formula}.

For a product $K=\prod F_j$, the additive character and self-dual
measure factor over $j$, as do the phase and the critical equations.
A compactly supported function invariant under a fixed product
lattice is a finite linear combination of products of characteristic
functions of balls in the factors.  Apply the field result to each
such product and use linearity and Fubini.  Enclose the projections
of $Q$ in compact subsets of $F_j^\times$ and take the product of the
resulting neighborhoods.  This proves the full assertion, with
uniformity for all the stated amplitudes.
\end{proof}

\begin{remark}\label{th-uniformity-limit}
The theorem gives a common neighborhood for the displayed fixed local
data.  It does not give a neighborhood independent of $m$, $p$, the
field, or the character conductor.  In particular the inverse Hessian
contains $(m+1)^{-1}$, and the analytic square root used in
Lemma~\ref{th-morse} depends on the valuation of $2$.  This dependence
cannot be discarded when passing to a varying global family.
The source's local theorem likewise quantifies its neighborhood after
the local field, character, dimension, amplitude scale, and compact
parameter set have been fixed.  Its later passage from these local
identities to a uniform positive majorant requires a separate argument.
\end{remark}

\subsection{Finite cokernels and positivity of the local kernel}

Let $R\subset K$ be either of the self-dual compact open orders above.
For $Z\in M_n(R)\cap\operatorname{GL}_n(K)$, define
\[
 C_Z=R^n/ZR^n,\qquad e_Z=[e_n],\qquad
 \lambda_Z([v])=e_n^tZ^{-1}v\pmod R\quad\text{in }K/R.
\]
For $n=0$ take $C_Z=0$, $e_Z=0$, and $\lambda_Z=0$.
The functional $\lambda_Z$ is well-defined because replacing $v$ by
$v+Zw$ adds $e_n^tw\in R$.  Call $(C_Z,e_Z,\lambda_Z)$ the marked
finite module associated to $Z$.  An isomorphism of two such modules
is an $R$-linear isomorphism carrying the marked element and the
functional to their counterparts.

\begin{lemma}\label{th-coker}
The module $C_Z$ is finite, and
\[
 \#C_Z=|\det Z|_K^{-1}.
\]
These assertions do not require $R$ to be a principal ideal domain.
\end{lemma}

\begin{proof}
The submodule $ZR^n$ is open and compact, since $Z$ is an invertible
continuous linear map.  It is contained in $R^n$.  The quotient is
therefore compact and discrete, hence finite.  All its cosets have
volume $\operatorname{vol}(ZR^n)=|\det Z|_K$, while
$\operatorname{vol}(R^n)=1$.
\end{proof}

Here is the precise counting mechanism for positivity.  Let $\mathcal W$
be a finite set.  For each $w\in\mathcal W$, choose a finite family
of marked finite $R$-modules $C_{w,q}$ and weights $a_{w,q}\in\mathbb C$.
Put
\begin{equation}\label{th-coker-kernel}
 \mathcal K(w,w')=
  \sum_{q,q'}a_{w,q}\overline{a_{w',q'}}
       \#\operatorname{Iso}(C_{w,q},C_{w',q'}).
\end{equation}
\begin{lemma}\label{th-coker-gram}
The matrix $\mathcal K$ is positive semidefinite.  More precisely, if
$C_\alpha$ runs through the isomorphism classes occurring, and
\[
 B_w(\alpha)=\sum_q a_{w,q}
                  \#\operatorname{Iso}(C_{w,q},C_\alpha),
\]
then
\begin{equation}\label{th-coker-gram-formula}
 \mathcal K(w,w')=
 \sum_\alpha
 \frac{B_w(\alpha)\overline{B_{w'}(\alpha)}}
      {\#\operatorname{Aut}(C_\alpha)}.
\end{equation}
\end{lemma}

\begin{proof}
If two modules are not isomorphic, both sides of the corresponding
identity for their isomorphism count are zero.  If they are isomorphic
to $C_\alpha$, each set of isomorphisms to $C_\alpha$ is a torsor
under the finite nonempty group $\operatorname{Aut}(C_\alpha)$.
Consequently
\[
 \#\operatorname{Iso}(C,D)=
 \sum_\alpha
 \frac{\#\operatorname{Iso}(C,C_\alpha)
       \#\operatorname{Iso}(D,C_\alpha)}
      {\#\operatorname{Aut}(C_\alpha)}.
\]
Substitution gives \eqref{th-coker-gram-formula}.  Its quadratic form
at coefficients $b_w$ is
$\sum_\alpha|\sum_w \overline{b_w} B_w(\alpha)|^2/
\#\operatorname{Aut}(C_\alpha)$, which is nonnegative.
\end{proof}

In the source, $w$ consists of positive block sizes $s_1,\ldots,s_l$
and scalars $t_1,\ldots,t_l\in K^\times$.  The matrix $Z(w)$ is block
diagonal with blocks $t_j$ times the reversing permutation matrix of
size $s_j$.  The index $q$ ranges over the integral matrices in its
specified lower- and upper-unipotent orbit, modulo integral
unipotent changes.  The weight is
$\psi(\sum x_{i+1,i}+\sum y_{i,i+1})$.
For every finite collection of $w$, the source proves that only
finitely many such cosets occur.  Once that finiteness and the
invariance under the integral changes have been established, its
local kernel is exactly \eqref{th-coker-kernel}.  Thus the positivity
step is Lemma~\ref{th-coker-gram}, with no appeal to a global
automorphic identity.

There is also an immediate support restriction.  All matrices in
the orbit of $Z(w)$ have determinant $\det Z(w)$, because the
unipotent factors have determinant one.  By Lemma~\ref{th-coker},
\begin{equation}\label{th-det-support}
 \mathcal K(w,w')\ne0\quad\Longrightarrow\quad
 |\det Z(w)|_K=|\det Z(w')|_K.
\end{equation}
The factor multiplying the kernel after adjoining a block of size
$s$ is
$|tt'|_K^{s(s-1)/2}$.  It is a product of a positive factor depending
on $t$ and one depending on $t'$, so preserves positive
semidefiniteness.  This establishes the positivity and determinant
support portions of that construction.  The stronger support
condition $t^{-1}-t'^{-1}\in R$ uses the additional averaging over
successive coordinate functionals; it does not follow merely from
\eqref{th-det-support}.

\subsection{Finite Fourier averaging and the local limit}

Two further steps of the local construction admit short proofs that
do not require an estimate for an oscillatory integral.  For Hermitian
matrices $M,N$, the notation $M\le N$ means that $N-M$ is positive
semidefinite.

\begin{lemma}\label{th-cut-contraction}
Let $G$ and $G_c$ be Hermitian matrices on a finite index set, with
$0\le G_c\le G$.  There are vectors $u_i$ in a finite-dimensional
Hilbert space and a linear contraction $T$ such that
\[
 G_{ij}=\langle u_i,u_j\rangle,\qquad
 (G_c)_{ij}=\langle Tu_i,Tu_j\rangle.
\]
\end{lemma}

\begin{proof}
Equip the vector space of coefficients with the positive semidefinite
form $a^*Gb$ and quotient by its null space.  The form $a^*G_cb$
descends to that quotient: a vector of zero $G$-norm has zero
$G_c$-norm, and the Cauchy--Schwarz inequality for $G_c$ makes its
pairing with every vector zero.  On the quotient the second form is
$\langle a,Bb\rangle$ for a self-adjoint operator $0\le B\le I$.
Take $T=B^{1/2}$ and let $u_i$ be the images of the coordinate vectors.
The spectral theorem gives $\|T\|\le1$ and both stated identities.
This also covers the zero-dimensional quotient.
\end{proof}

\begin{proposition}\label{th-fourier-average}
Let $A$ be a finite nonempty set, let $B$ be a finite abelian group,
and put $a=\#A$, $b=\#B$.  Let $I$ be a finite index set and
$u_{\rho,t,i}$ vectors in a Hilbert space $\mathcal E$, indexed by
$\rho\in A$, $t\in B$, and $i\in I$.  Suppose that
\[
 \langle u_{\rho,t,i},u_{\rho',t',j}\rangle=0
 \quad\text{if }\rho\ne\rho'.
\]
For each character $\chi\in\widehat B$ choose numbers
$\omega_\rho(\chi)$ of modulus one, and put
\[
 v_{\rho,\chi,i}=\frac1b\sum_{t\in B}
             \overline{\chi(t)}u_{\rho,t,i},\qquad
 v_{\chi,i}=\frac1a\sum_{\rho\in A}
             \omega_\rho(\chi)v_{\rho,\chi,i},
\]
\[
 \Theta_\chi(i,j)=\frac1a\sum_{\rho\in A}
          \langle v_{\rho,\chi,i},v_{\rho,\chi,j}\rangle.
\]
If $T:\mathcal E\to\mathcal E$ is a linear contraction and
$N_\chi(i,j)=a\langle Tv_{\chi,i},Tv_{\chi,j}\rangle$, then
\begin{equation}\label{th-fourier-domination}
 0\le N_\chi\le\Theta_\chi,
 \qquad
 \sum_{\chi\in\widehat B}\Theta_\chi(i,j)
 =\frac1{ab}\sum_{\rho\in A}\sum_{t\in B}
       \langle u_{\rho,t,i},u_{\rho,t,j}\rangle.
\end{equation}
Consequently, if the matrix on the right is at most a fixed matrix
$M$, then $\sum_\chi\Theta_\chi\le M$.
\end{proposition}

\begin{proof}
For coefficients $c_i$, write
$v_{\rho,\chi,c}=\sum_i c_i v_{\rho,\chi,i}$ and similarly
for $v_{\chi,c}$.  Orthogonality between distinct $\rho$ gives
\[
 \sum_{i,j}\overline{c_i}c_j\Theta_\chi(i,j)
 =\frac1a\sum_\rho\|v_{\rho,\chi,c}\|^2
 =a\|v_{\chi,c}\|^2
 \ge a\|Tv_{\chi,c}\|^2.
\]
This proves the first assertion, including positivity.  Character
orthogonality gives
\[
 \sum_{\chi\in\widehat B}\chi(t)\overline{\chi(t')}
 =b\,1_{t=t'}.
\]
This follows either by decomposing $B$ into cyclic
groups and summing geometric series, or by translating this
character sum.  Expanding the inner products and applying that
identity proves the second assertion.  The last assertion follows
directly by comparison of quadratic forms.
\end{proof}

Here is the exact role of this proposition in the local induction.
One chooses additive lattices
$L_1\subset R\subset L_2$ and representatives $A$ of $L_2/R$
and $B$ of $R/L_1$.  The latter quotient is the group in
Proposition~\ref{th-fourier-average}; representatives of $A$ are
only a finite indexing set.  Each shift $\sigma=\rho+t$ changes
the small parameter $\tau$ to $\tau_\sigma$ by
$\tau_\sigma^{-1}=\tau^{-1}+\sigma$.
The source's unrestricted orbital sum supplies a Gram matrix on the
indices $(\rho,t,i)$.  Restricting the orbital sum to a subset
supplies a smaller positive semidefinite matrix, so
Lemma~\ref{th-cut-contraction} supplies $T$.
The support condition in the preceding paragraph gives the required
orthogonality when $\rho\ne\rho'$.
Self-duality identifies $\widehat{R/L_1}$ with $L_1^\perp/R$.
The phases $\omega_\rho$ account for the representatives of
$L_2/R$.  The crucial remaining identity says that the matrix
$a\langle Tv_{\chi,i},Tv_{\chi,j}\rangle$ equals the desired
kernel after adjoining the larger block, for parameters in the
chosen compact annulus.  Once that identity and the bound for the
averaged unrestricted diagonal have been established,
\eqref{th-fourier-domination} proves both the positive majorant and
its total-mass bound.  Thus the averaging inequality itself adds
neither an unspecified constant nor a loss depending on $a$ or $b$.

The following compactness argument removes the annulus.  A function
on $K$ is called $R$-periodic if it is constant on every additive
coset of $R$.

\begin{proposition}\label{th-local-limit}
Let $L_c\subset K$ be a compact open additive subgroup containing
$R$, and let $I$ be a finite index set.  Let $M$ be a fixed positive
semidefinite matrix on $I$.  Suppose that $N(t)$ is a Hermitian
matrix for every $t\in K^\times$, with $N(t)=0$ outside $L_c$.
Let $\Omega_j$ be increasing subsets of $K^\times$ whose union is
$K^\times$.  Suppose that for every $j$ there is an $R$-periodic
matrix-valued function $\Theta_j$ such that
\[
 \Theta_j(x)\ge0,\qquad
 \Theta_j(x)=0\ (x\notin L_c),\qquad
 \sum_{x\in K/R}\Theta_j(x)\le M,
\]
and $\Theta_j(t)\ge N(t)$ for $t\in\Omega_j$.
Then there is an $R$-periodic function $\Theta$ with the same
positivity, support, and total-mass properties and
$\Theta(t)\ge N(t)$ for every $t\in K^\times$.
Any fixed collection of zero matrix entries shared by all
$\Theta_j(x)$ is shared by $\Theta(x)$.  If $L_c=R$, then
\begin{equation}\label{th-unit-head-bound}
 0\le\Theta(x)\le1_R(x)M\qquad(x\in K).
\end{equation}
\end{proposition}

\begin{proof}
The quotient $L_c/R$ is compact and discrete, hence finite.  Thus
the matrices $\Theta_j$ are described by a fixed finite number of
complex entries.  For every coset $x$ one has
$0\le\Theta_j(x)\le M$, because all other summands are positive
semidefinite.  The $2$ by $2$ principal minors therefore give
\[
 |\Theta_j(x)_{il}|^2
 \le\Theta_j(x)_{ii}\Theta_j(x)_{ll}
 \le M_{ii}M_{ll}.
\]
A subsequence converges in this finite-dimensional space.  Define
$\Theta$ by its limits on these cosets and set it to zero outside
$L_c$.  The cone of positive semidefinite matrices is closed:
for every vector $c$, the inequality $c^*Ac\ge0$ passes to an
entrywise limit.  This proves positivity and the total-mass bound,
since the latter is a finite sum.  For fixed $t\in K^\times$,
one has $t\in\Omega_j$ for all sufficiently large $j$, so the
same closedness argument gives $\Theta(t)\ge N(t)$.
Zero entries also pass to the limit.  Finally, if $L_c=R$, the
sum over $K/R$ has just one possibly nonzero term, proving
\eqref{th-unit-head-bound}.
\end{proof}

For the present product of local fields one may take $\Omega_j$ to
be the product of the annuli on which each component valuation lies
between $-j$ and $j$.  These sets are compact, open, exhaust the
units of $K$, and are stable under multiplication by any root of
unity.  If the source's additional support assertion
\[
 N(t)_{hh'}\ne0\quad\Longrightarrow\quad
       (\det Z(h))t\in R
\]
has been proved, a common permissible subgroup is the compact open
$R$-module
$L_c=R+\sum_h(\det Z(h))^{-1}R$, where $h$ ranges over the fixed
finite family of previous block configurations and $Z(h)$ is their
block reversal matrix defined above.  If all $Z(h)$ and their inverses have
entries in $R$, then each determinant is a unit of $R$, so
$L_c=R$.  In this case the upper bound in
\eqref{th-unit-head-bound} is precisely the final extra condition
in the local majorant statement, after substituting its matrix
$M$.  No rate of convergence, uniform stationary-phase radius,
or infinite-dimensional compactness theorem is needed for this
last passage.  The subsequent local reconstruction proves that
the truncated orbital Gram matrix equals the integral of the
larger-block kernel.  It establishes the actual orbital germ
formula and the cancellation between distinct stationary points.

\subsection{Reproducing constants and convolution on a quotient}

The possible scalar in the local reproducing formula can be determined
without evaluating a quotient integral point by point.  We state the
needed functional-analytic argument explicitly.

Let $G$ be a second countable unimodular locally compact group with a
fixed Haar measure.  Let $\pi$ be a unitary representation on a Hilbert
space $V$, and suppose $d_\pi>0$ satisfies the formal-degree identity
\begin{equation}\label{th-schur}
 \int_G\langle a,\pi(g)b\rangle
       \overline{\langle a',\pi(g)b'\rangle}\,dg
   =d_\pi^{-1}\langle a,a'\rangle\langle b',b\rangle
\end{equation}
for all vectors.  The existence of this identity for an irreducible
discrete series is a representation-theoretic input; it is not being
inferred from unitarity alone.  Choose $w\in V$ of norm one and put
\[
 p(g)=d_\pi\overline{\langle w,\pi(g)w\rangle}.
\]
In addition to \eqref{th-schur}, assume $p\in L^1(G)$.

\begin{lemma}\label{th-projection}
For convolution $(p*q)(x)=\int_Gp(g)q(g^{-1}x)\,dg$ one has
$p*p=p$ and $p^*(g)=\overline{p(g^{-1})}=p(g)$.
In any strongly continuous unitary representation $R$ of $G$, the
bounded operator
\[
 R(p)v=\int_Gp(g)R(g)v\,dg
\]
is an orthogonal projection.  In particular, if $v\ne0$ and
$R(p)v=cv$ with $c\ne0$, then $c=1$.
\end{lemma}

\begin{proof}
The coefficient belongs to $L^2$ by \eqref{th-schur}, so the
convolution integrand is absolutely integrable by Cauchy--Schwarz.
Unitarity and \eqref{th-schur} give
\begin{align*}
 (p*p)(x)
 &=d_\pi^2\int_G
    \overline{\langle w,\pi(g)w\rangle}
    \langle\pi(x)w,\pi(g)w\rangle\,dg\\
 &=d_\pi\langle\pi(x)w,w\rangle=p(x).
\end{align*}
The identity $p^*=p$ follows directly from unitarity.
The integral defining $R(p)v$ exists in norm and
$\|R(p)\|\le\|p\|_1$.  Fubini is justified by the majorant
$|p(g)q(h)|\|v\|$, and gives $R(p)R(q)=R(p*q)$.
Invariance under inversion gives $R(p)^*=R(p^*)$.
Thus $R(p)$ is self-adjoint and idempotent, hence an orthogonal
projection.  Applying idempotence to $v$ gives $c^2v=cv$;
$c\ne0$ and $v\ne0$ imply $c=1$.
\end{proof}

For application to the source, let $S=N\rtimes H$ be the local
stabilizer, with $N$ unipotent and $H$ compact.  Its action on $N$
preserves Haar measure: its modular factor is a continuous
homomorphism from the compact group $H$ to $\mathbb R_{>0}$ and
therefore is one.  Thus $S$ is unimodular.  Let $\rho$ be its strongly
continuous unitary representation, formed from the specified character of $N$ and
the compact representation of $H$.  The induced Hilbert space is
\[
 \mathcal I=\left\{F:G\to V_\rho:
 F(sg)=\rho(s)F(g),\quad
 \int_{S\backslash G}\|F(g)\|^2\,d\bar g<\infty\right\},
\]
where the functions are measurable and functions equal almost
everywhere are identified.  Right
translation $R(g)F(x)=F(xg)$ is unitary, by the quotient integration
formula.  Lemma~\ref{th-projection} applies to this representation.
It proves that the nonzero scalar in a reproducing formula for a
nonzero $F\in\mathcal I$ must be one.

The nonvanishing, membership in $\mathcal I$, and the $L^1$ condition
are substantive hypotheses.  The last does not follow from square
integrability of discrete-series coefficients.  In the source it is
obtained from a weighted coefficient estimate under a large
Harish--Chandra parameter gap.  The argument above replaces the
subsequent scalar-normalization argument, not that weighted estimate.

\subsection{Unfolding the Gram matrix}

Let $\Gamma\subset G$ be a discrete subgroup, and let
$\Gamma_u\subset\Gamma$ be subgroups.  Let $C_u$ be finite-dimensional
Hilbert spaces and let $\phi_{u,v}(g)$ depend linearly on $v\in C_u$,
with $\phi_{u,v}$ left $\Gamma_u$-invariant.  Define
\[
 P_uv(g)=\sum_{\delta\in\Gamma_u\backslash\Gamma}
                         \phi_{u,v}(\delta g).
\]
Assume the series converge absolutely almost everywhere, give
$L^2(\Gamma\backslash G)$ functions, and that for every pair of
vectors from chosen bases,
\begin{equation}\label{th-unfold-integrable}
 \int_{\Gamma_u\backslash G}|\phi_{u,v}(g)|
     \sum_{\delta\in\Gamma_{u'}\backslash\Gamma}
                  |\phi_{u',w}(\delta g)|\,d\bar g<\infty.
\end{equation}

\begin{lemma}\label{th-unfold}
Choose representatives $\eta$ for
$\Gamma_{u'}\backslash\Gamma/\Gamma_u$, and put
$J_\eta=\Gamma_u\cap\eta^{-1}\Gamma_{u'}\eta$.
With compatible quotient measures,
\begin{equation}\label{th-unfold-formula}
 \langle P_uv,P_{u'}w\rangle
 =\sum_\eta\int_{J_\eta\backslash G}
       \overline{\phi_{u,v}(g)}\phi_{u',w}(\eta g)\,d\bar g.
\end{equation}
The right side is absolutely convergent.
\end{lemma}

\begin{proof}
Expand the first Poincar\'e series and integrate on a measurable
fundamental domain for $\Gamma$.  Translating that domain by the
representatives of $\Gamma_u\backslash\Gamma$ tiles a fundamental
domain for $\Gamma_u$.  Since $P_{u'}w$ is left $\Gamma$-invariant,
the result is
$\int_{\Gamma_u\backslash G}
\overline{\phi_{u,v}(g)}P_{u'}w(g)\,d\bar g$.
For each double coset $\Gamma_{u'}\eta\Gamma_u$, its left cosets
are $\Gamma_{u'}\eta\alpha$, with
$\alpha\in J_\eta\backslash\Gamma_u$.
Expanding this part of the second series and unfolding once more
gives the integral over $J_\eta\backslash G$ in
\eqref{th-unfold-formula}.  Tonelli applied to the absolute values,
using \eqref{th-unfold-integrable}, justifies both rearrangements
and proves the asserted convergence.  Finite-dimensionality extends
the basis assumption to all vectors.
\end{proof}

By \eqref{th-normalized-K}, division of each term by $\sqrt{c_uc_{u'}}$
gives the corresponding expansion of $K_{uu'}$.  This explains which
kernel is being bounded: it is the Gram operator of the actual
Poincar\'e maps, rather than an unrelated matrix with a similar name.
Further decomposition of the double cosets into flags, lattices,
and local orbits, and the estimates for their absolute sums, remain
arithmetic work.  The hypotheses in Lemma~\ref{th-unfold} must be
proved for the proposed seeds; they are not consequences of the
formal definition of a Poincar\'e series.

\subsection{From the orbital estimates to the pairwise bound}

The last combination of the orbital estimates admits a short proof.
We give it because it specifies exactly what remains to be estimated.
Let $a_t$ be an absolutely summable family in a Hilbert space
$\mathcal A$, and let $v_t$ be an absolutely summable family in a
Hilbert space $\mathcal V$.  Let $r(t)$ be a row label, and let
$i(r)$ belong to a finite set $I$.  Suppose that for each of two
disjoint sets of terms $S=\mathcal G,\mathcal N$,
\begin{equation}\label{th-row-norm}
 \left\|\sum_{t\in S}a_t\right\|^2
  =\operatorname{Re}\sum_{i,j\in I}\Gamma_{ij}
       \left\langle\sum_{t\in S:i(r(t))=i}v_t,
                       \sum_{t\in S:i(r(t))=j}v_t\right\rangle.
\end{equation}
Assume $P\gamma I-\Gamma$ is positive semidefinite, with
$P\ge0$, $\gamma>0$, and assume
\begin{equation}\label{th-row-orth}
 \langle v_t,v_{t'}\rangle=0
 \quad\text{if }i(r(t))=i(r(t'))\text{ and }r(t)\ne r(t').
\end{equation}
In the application, $\mathcal A$ is a space of Hilbert--Schmidt
operators.  The sets $\mathcal G$ and $\mathcal N$ consist of the
generic and nongeneric terms with nonzero Gram determinant.

For equal-row pairs of nonzero-determinant terms, suppose an integer
$s(t,t')\ge1$ and a scalar/non-scalar distinction at $s=1$ are given.
For a set $\mathcal Q$ of ordered pairs define
\[
 E(\mathcal Q)=\gamma\sum_{(t,t')\in\mathcal Q}
                    \operatorname{Re}\langle v_t,v_{t'}\rangle,
 \qquad
 E_{\mathrm{abs}}(\mathcal Q)=\gamma\sum_{(t,t')\in\mathcal Q}
                    |\langle v_t,v_{t'}\rangle|.
\]
All these sums converge absolutely.  Put $\mathcal B$ for the
equal-row pairs with $s\ge2$, $\mathcal C$ for the generic
equal-row pairs with $s=1$ that are not scalar, $\mathcal M$ for
the scalar generic pairs with $s=1$, and $\mathcal U$ for the
nongeneric pairs with $s=1$.

\begin{proposition}\label{th-pair-bound}
Let $d\ge0$ and $A,\epsilon,\epsilon'\ge0$.  Suppose
\begin{equation}\label{th-four-estimates}
 E_{\mathrm{abs}}(\mathcal B)\le\epsilon d,\quad
 E_{\mathrm{abs}}(\mathcal C)\le\epsilon d,\quad
 E(\mathcal M)\le A d,\quad E(\mathcal U)\le\epsilon'd.
\end{equation}
Then
\begin{equation}\label{th-pair-output}
 \left\|\sum_{t\in\mathcal G}a_t+
             \sum_{t\in\mathcal N}a_t\right\|^2
 \le P\left(\sqrt{A+2\epsilon}
                  +\sqrt{\epsilon+\epsilon'}\right)^2d.
\end{equation}
\end{proposition}

\begin{proof}
The positive semidefinite inequality and
\eqref{th-row-norm} bound the squared norm for a set $S$ by
\[
 P\gamma\sum_{i\in I}
       \left\|\sum_{t\in S:i(r(t))=i}v_t\right\|^2.
\]
Expand the squares.  Equation \eqref{th-row-orth} removes pairs
with different row labels.  For $S=\mathcal G$, the remaining
pairs partition into $s\ge2$, scalar $s=1$, and nonscalar $s=1$;
for $S=\mathcal N$, they partition into $s\ge2$ and $s=1$.
Using \eqref{th-four-estimates} gives bounds
$P(A+2\epsilon)d$ and $P(\epsilon+\epsilon')d$ respectively.
The triangle inequality gives \eqref{th-pair-output}.  Absolute
summability justifies every expansion and restriction of the pair sums.
\end{proof}

\begin{lemma}\label{th-small-errors}
If $b\ge16$ and
\[
 0\le P\le e^{6b^{-7}},\quad
 0\le A\le(1+b^{-7})e^{b^{-14}},\quad
 0\le\epsilon,\epsilon'\le b^{-14},
\]
then the multiplier in \eqref{th-pair-output} is at most
\begin{equation}\label{th-sharper-error}
 1+14b^{-7}\le1+b^{-6}.
\end{equation}
\end{lemma}

\begin{proof}
Write $u=b^{-7}\le1/1000$.  The elementary bound
$e^x\le1+x+2x^2$ for $0\le x\le1/2$ follows from
$e^x(1-x)\le1$ and $1/(1-x)\le1+x+2x^2$.
It gives $P\le1+7u$ and $e^{u^2}\le1+2u^2$.
Consequently
\[
 A+2\epsilon\le1+u+5u^2\le1+2u,
 \qquad P(A+2\epsilon)\le1+10u.
\]
Also $P(\epsilon+\epsilon')\le\tfrac94u^2$.
The sum of the two square roots is therefore at most
$1+5u+\tfrac32u=1+\tfrac{13}{2}u$.  Squaring and using
$u\le1/1000$ bounds it by $1+14u$.
Finally $14u\le bu=b^{-6}$.
\end{proof}

The slightly stronger bound $14b^{-7}$ is available without changing
the orbital estimates.  It improves the last numerical combination;
it does not eliminate any of the four estimates in
\eqref{th-four-estimates}.

Here is a second useful simplification of the numerical part.  Suppose
local quantities $q_{\mathfrak p}$ satisfy
$q_{\mathfrak p}\ge p$ at a place above the rational prime $p$, and
there are at most $f$ such places above each $p$.  For $a>1$,
\begin{equation}\label{th-density-bound}
 \prod_{\mathfrak p}\exp(4q_{\mathfrak p}^{-a})
 \le\exp\left(4f\,2^{-a}\left(1+\frac2{a-1}\right)\right).
\end{equation}
Indeed
$\sum_{\mathfrak p}q_{\mathfrak p}^{-a}
\le f\sum_{n\ge2}n^{-a}
\le f(2^{-a}+\int_2^\infty t^{-a}\,dt)$.
This also proves convergence, and applies to every finite subproduct.
The source uses $a=k-2r-2\ge2$ and the weaker upper bound
$\exp(12f2^{-a})$.  Thus its local-density accumulation does not
conceal a constant depending on the number of places.

Likewise, for $q\ge3$,
\[
 \prod_{j=1}^k(1-q^{-j})^{-1}
 \le\exp\left(\frac{q}{(q-1)^2}\right)
 \le\exp\left(\frac9{4q}\right).
\]
Use $-\log(1-z)\le z/(1-z)$ and sum the geometric series.
If two split primes have sizes $g_1,g_2\ge(kb)^8$ and each has $f$
places, the product of these bounds is at most
$\exp(9f/(2(kb)^8))$.  This is smaller than the coarser
$\exp(6b^{-7})$ used in Lemma~\ref{th-small-errors} when
$f\le k$, $k\ge1$, and $b\ge1$.
For clarity, the actual prefactor $P$ in the source's row estimate
is the different product
\[
 P=\prod_{v\mid g_1g_2}(1+3/q_v)
   =\bigl((1+3/g_1)(1+3/g_2)\bigr)^f.
\]
The equality uses complete splitting at these two rational primes.
It satisfies the required bound directly:
\[
 P\le\exp(3f/g_1+3f/g_2)
   \le\exp\bigl(6f/(kb)^8\bigr)
   \le\exp(6b^{-7}).
\]
Thus the separate Euler-product estimate above is not needed to
identify this prefactor.

\subsection{The remaining global statement, with its scale visible}

We first specify all parameters in the following conditional numerical
combination.  The integers $k,f\ge1$ and $r\ge0$ have the following
meaning in the source.  The CM extension is $E/F$, with $F$ totally
real of degree $f$ over $\mathbb Q$.  The relevant Hermitian space has
dimension $k+2r+1$ over $E$ and signature $(k+r,r+1)$ at each real
place of $F$.  The integer $r$ is the length of the isotropic flag,
and the compact stabilizer has dimension $k$ in its defining
representation.  The number of successive lifts is $2r+2$.  The
positive real number $b$ is the height scale for the parameters of
the representations; in the growing family it is $\exp(x^\varepsilon)$,
where $x>1$ tends to infinity and $\varepsilon>0$ is fixed.

The real numbers $c_5>0$ and $C_n\ge0$ are respectively the decay
constant in the large-corank and nonscalar bounds and the growth
constant in the nongeneric bound.  The concrete source takes
$c_5=1/3$ and chooses $C_n$ after the fixed auxiliary constants,
independently of the subsequent growing parameters.  The positive
real number $C_\Gamma$ is chosen after the constants in the diagonal
estimates, and
\[
 \Gamma=\exp\bigl(C_\Gamma(2r+2)^5f\log b\bigr).
\]
The real number $C_0\ge1$ measures the assumed comparison between
the dimensions $d_u$ and $D_*/\Gamma^2$: the source requires
$C_0^{-1}\le d_u\Gamma^2/D_*\le C_0$ for every $u$.
Only its upper bound enters the diagonal estimate below.

Here $D_*>0$ is the proposed main term for the dimension of the
ambient automorphic space.  More precisely, write $G$ for the unitary
group just specified, $d_\infty>0$ for the product of the
archimedean formal degrees with the chosen Haar measures, and
$e_{\mathrm{fin}}$ for the finite-adelic convolution idempotent
specifying the finite-level condition.  Its normalization is
\[
 D_*=d_\infty\,
 \operatorname{vol}\bigl(G(F)\backslash G(\mathbb A_F)\bigr)
 \operatorname{Re} e_{\mathrm{fin}}(1).
\]
Both the volume and the idempotent use the measures fixed in the
automorphic construction.  This definition does not itself assert
an asymptotic formula for the actual dimension.  In the numerical
argument that follows, $k,f,r,b,c_5,C_n,C_\Gamma,C_0,D_*$ are explicit
input numbers with the domains just stated and the displayed
inequalities; their realization by the particular global family
remains a separate assertion.

For distinct indices $u,u'$, put
$d=d_ud_{u'}/D_*$, where $d_u=\dim C_u$ and $D_*>0$.
The actual four estimates in the source have the following scales:
\begin{align}\label{th-global-scales}
 \epsilon&=\exp(-c_5kf\log k),\notag\\
 \epsilon'&=2^{-(k-2r-5)}
       \exp\bigl(C_n(2r+2)^5f\log b\bigr),\notag\\
 A&=(1+\delta)e_1,\qquad
 0\le\delta\le b^{-7},\qquad
 1\le e_1\le\exp(12f2^{-(k-2r-2)}).
\end{align}
The dimensional restrictions include
\begin{equation}\label{th-global-parameters}
 \begin{gathered}
 b\ge16,\quad k\ge32,\quad f\le k,\quad 4r+12\le k,
 \quad70\log b\le k,\\
 5c_5f\log k\ge1,\qquad
 8C_n(2r+2)^5f\log b\le k.
 \end{gathered}
\end{equation}
They imply $\epsilon,\epsilon'\le e^{-k/5}\le b^{-14}$ and
$12f2^{-(k-2r-2)}\le e^{-k/5}$.  To verify the latter bound, use
$k-2r-2\ge k/2+4$, so that the first expression is at most
$f2^{-k/2}\le k2^{-k/2}\le e^{-k/5}$ for $k\ge32$.
The last elementary inequality follows because
$\log k\le k((\log2)/2-1/5)$ at $k=32$ and the right side has
larger derivative thereafter.  For $\epsilon'$, use
$k-2r-5\ge k/2$ and
$C_n(2r+2)^5f\log b\le k/8$; then
$-(k/2)\log2+k/8\le-k/5$.

Assuming the actual orbital decomposition satisfies
\eqref{th-row-norm}--\eqref{th-four-estimates}, and assuming the
degenerate and identity terms vanish for $u\ne u'$, the already
proved proposition and lemma give
\begin{equation}\label{th-off-final}
 \|K_{uu'}\|_{\mathrm{HS}}^2
 \le(1+14b^{-7})\frac{d_ud_{u'}}{D_*}.
\end{equation}
The source states the weaker $(1+b^{-6})$ version.

On the diagonal write $K_{uu}=I+Z_u+R_u$, where $Z_u$ is the sum
of the nonidentity terms with zero Gram determinant.  The separate
estimates required there are
\begin{equation}\label{th-diag-input}
 \|Z_u\|_{\mathrm{HS}}^2\le d_u e^{-kf\log k},\qquad
 \|R_u\|_{\mathrm{HS}}^2\le\Gamma\frac{d_u^2}{D_*},
 \qquad \frac{d_u\Gamma^2}{D_*}\le C_0.
\end{equation}
They imply the slightly sharper diagonal defect
\begin{equation}\label{th-diag-final}
 \|K_{uu}-I\|_{\mathrm{HS}}^2
 \le d_u\left(e^{-kf\log k/2}+\sqrt{C_0/\Gamma}\right)^2.
\end{equation}
This is just the triangle inequality together with the last
inequality in \eqref{th-diag-input}.  All constants and their roles
are explicit at this stage.

The Hilbert--Schmidt norm of a matrix representing a linear map in
orthonormal bases is the operator's Hilbert--Schmidt norm.
Consequently, if $M_u$ is the orthonormal matrix of $P_u$ and
$A_u=c_u^{-1/2}M_u$, then
\[
 A_u^*A_{u'}=K_{uu'}.
\]
Equations \eqref{th-off-final} and \eqref{th-diag-final} therefore
give exactly the pairwise and diagonal estimates needed by the
previous slope comparison.  Compression by an isometric inclusion
of the ramified subspace is handled by the compression lemma proved
there.  This last identification requires no additional estimate.

\paragraph{Local statements established in the following sections.}
The local reconstruction proves the
finiteness and representative independence of the actual local
sums, their local constancy uniformly against all comparison
configurations, and the strong inverse-parameter support.  It also
proves the root-fiber change of variables and the cancellation of
terms from distinct stationary points.  The fixed-initial-block
construction and the Hessenberg change of variables establish the
actual kernel germ in Theorem~\ref{ga-actual-germ}, hence the exact
truncated orbital identity.  The finite character identity and its
averaged diagonal bound then complete the positive-majorant
induction in Theorem~\ref{tp-majorant}.
Propositions~\ref{th-fourier-average} and \ref{th-local-limit}
supply its averaging and limiting steps with the asserted constants.

\paragraph{What has not yet been proved in this reconstruction.}
The following are concrete remaining assertions about the proposed
automorphic family, rather than consequences of the local lemmas.
First, the actual seeds must have the weighted integrability and
nonvanishing properties used above; the source obtains these from
discrete-series and Bessel-model theorems with large parameter gaps.
Second, the global recursion after Poisson summation must
bound the scalar generic pair sum by $(1+\delta)e_1d$ with
$\delta\le b^{-7}$.  Its arithmetic inputs include the lattice
indices and the sums over the dual lattices, with constants uniform
in the stated growing parameters.  Third, the larger-corank,
nonscalar, and nongeneric sums in \eqref{th-global-scales}, the
degenerate diagonal bound in \eqref{th-diag-input}, and the required
ambient-dimension estimate must be established for these particular
maps.  Positivity and stationary phase alone imply none of those
global estimates.

No counterexample satisfying all the actual hypotheses of the local
results proved here has been found.  This is not an assertion of an
exhaustive audit of the theta energy recursion.  The source trace
accompanying this section distinguishes the statements and proofs
read in full from the still unresolved arithmetic steps.

\subsection{Exact small controls and excluded shortcuts}

The accompanying program uses only integer arithmetic and reduction
modulo cyclotomic polynomials.  For $a=p^e$, $z=1$, and
$L=\mathbb Z_p^m$, it checks the finite sums
\[
 G_{p,m,e}=\sum_{x\in(\mathbb Z/p^e\mathbb Z)^m}
 \exp\left(\frac{2\pi i}{p^e}
       \left(\sum_i x_i^2+\sum_{i<j}x_ix_j\right)\right).
\]
The Gaussian lemma predicts
$|G_{p,m,e}|^2=p^{em+v_p(m+1)}$ when the admissibility condition
holds.  The checked cases are
$(p,m,e)=(2,1,2),(2,3,3),(3,2,2),(3,3,1),(5,2,1),(5,4,2)$.
The checks include residue characteristics dividing both $2$ and
$m+1$.  It also verifies the full stationary phase identity for
$m=1$, $D=1$, and amplitude $1_{\mathbb Z_p^\times}$ at
$(p,e)=(3,2),(3,3),(5,2),(2,8)$, using critical balls of radius
$p^{-1}$ for odd $p$ and $2^{-3}$ for $p=2$.

Two exact failures explain hypotheses that must be retained.
For $(p,m,e)=(2,3,1)$ one has $G_{2,3,1}=-4$, so the normalized
squared integral is $1/4$, whereas the determinant expression
without admissibility would be $1/2$.  The dual quadratic form is
not trivial on $L^\perp$ in this example.  This is a counterexample
to dropping \eqref{th-gauss-admissible}, not to
Lemma~\ref{th-gaussian} or to the source's corresponding statement.
For the split order \eqref{th-split-order}, with $F=\mathbb Q_p$
and odd $p$, the integral of $\psi_R(x^2,y^2)$ over
$\mathbb Z_p^2$ is $1/p$ for product Haar measure and $1$ for the
correct self-dual measure.  Thus the order's factor $p$ cannot be
omitted.  Both values are checked exactly for $p=3,5$.

Finally, a degenerate Hessian really is excluded by characteristic
zero, not by a condition on residue characteristic.  Over a field
of characteristic $p$, the same phase at $m=p-1$ has Hessian
determinant zero at its critical point.  The Gaussian argument then
does not apply.  The actual source assumes that every positive
integer is invertible in the local algebra, so this example does
not satisfy its hypotheses.



